\section{Methods}
\label{sec:methods}

\subsection{Unified Manifold Workspace and Coordinate Mapping}
\label{sec:methods:manifold}

\IfFileExists{sections/generated/params.tex}{% Auto-generated by scripts/TexExporter.jl. Do not edit manually.
% Generated UTC: 2026-06-12T18:41:40
\providecommand{\DzMaxMeters}{6.514}
\renewcommand{\DzMaxMeters}{6.514}
\providecommand{\DzMinMeters}{0.009}
\renewcommand{\DzMinMeters}{0.009}
\providecommand{\EnergyFloorThreshold}{1.0e-04}
\renewcommand{\EnergyFloorThreshold}{1.0e-04}
\providecommand{\GridNodes}{33}
\renewcommand{\GridNodes}{33}
\providecommand{\KappaMassBaseline}{1061}
\renewcommand{\KappaMassBaseline}{1061}
\providecommand{\ProfilesCampaignPool}{0}
\renewcommand{\ProfilesCampaignPool}{0}
\providecommand{\ProfilesStableWindow}{0}
\renewcommand{\ProfilesStableWindow}{0}
\providecommand{\ProfilesTotal}{0\xspace}
\renewcommand{\ProfilesTotal}{0\xspace}
\providecommand{\RankCompression}{0}
\renewcommand{\RankCompression}{0}
\providecommand{\RankMedian}{0.00}
\renewcommand{\RankMedian}{0.00}
\providecommand{\RankStd}{0.00}
\renewcommand{\RankStd}{0.00}
\providecommand{\SilhouettePeak}{0.000}
\renewcommand{\SilhouettePeak}{0.000}
\providecommand{\SpectralOrder}{32}
\renewcommand{\SpectralOrder}{32}
\providecommand{\StretchAlpha}{2.5}
\renewcommand{\StretchAlpha}{2.5}
\providecommand{\SvdToleranceFloor}{7.33e-15}
\renewcommand{\SvdToleranceFloor}{7.33e-15}
\providecommand{\TowerMax}{55.00\,\mathrm{m}}
\renewcommand{\TowerMax}{55.00\,\mathrm{m}}
\providecommand{\TowerMaxMeters}{55.00}
\renewcommand{\TowerMaxMeters}{55.00}
\providecommand{\TowerMin}{1.50\,\mathrm{m}}
\renewcommand{\TowerMin}{1.50\,\mathrm{m}}
\providecommand{\TowerMinMeters}{1.50}
\renewcommand{\TowerMinMeters}{1.50}
\providecommand{\TowerObservationLevels}{8}
\renewcommand{\TowerObservationLevels}{8}
\providecommand{\WindowSEarly}{0.000}
\renewcommand{\WindowSEarly}{0.000}
\providecommand{\WindowSIOP}{0.000}
\renewcommand{\WindowSIOP}{0.000}
\providecommand{\WindowSTransitional}{0.000}
\renewcommand{\WindowSTransitional}{0.000}
\providecommand{\alphaStretchBaseline}{2.5}
\renewcommand{\alphaStretchBaseline}{2.5}
}{}

Standard pseudospectral frameworks relying on uniform or Chebyshev-cosine grids distribute computational degrees of freedom symmetrically across the spatial domain. Such configurations are fundamentally unoptimized for the acute vertical gradients characterizing the nocturnal stable boundary layer (SBL), where sub-mesoscale gravity waves and intense shear instabilities concentrate within the lowest meters of the surface layer.

To preferentially focus node density where structural geometric derivatives are maximized, we establish a \emph{UnifiedManifoldWorkspace} utilizing an analytical hyperbolic tangent coordinate stretching transformation, $\mathcal{T}(\xi) \rightarrow z$. This operator transforms the bounded, isotropic computational coordinate $\xi \in [-1, 1]$ into the physical elevation domain defined by the CASES-99 tower instruments ($z_{\min} \le z \le z_{\max}$):
\begin{equation}
z(\xi) = z_c + \frac{L_d}{2} \left[ \frac{\tanh(\alpha_s \xi)}{\tanh(\alpha_s)} \right]
\label{eq:tanh_stretch}
\end{equation}
where $L_d = z_{\max} - z_{\min}$ represents the physical domain span, $z_c = (z_{\max} + z_{\min})/2$, and $\alpha_s \in (0, \infty)$ denotes the non-dimensional grid compression parameter. Differentiating Eq.~(\ref{eq:tanh_stretch}) with respect to the computational coordinate yields the analytical transformation metric Jacobian, $J(\xi)$:
\begin{equation}
J(\xi) = \frac{dz}{d\xi} = \frac{L_d}{2} \left[ \frac{\alpha_s}{\tanh(\alpha_s)} \right] \operatorname{sech}^2(\alpha_s \xi)
\label{eq:jacobian_stretched}
\end{equation}

Fixing the operational baseline at $\alpha_s = 2.50$ over a compressed physical domain span compresses the local node spacing to an ultra-dense scale directly above the surface boundary ($z_{\min} = 1.5\,\text{m}$), successfully clustering a significant fraction of the $N+1$ collocation nodes within the low-level jet (LLJ) shear core. To justify this spatial geometry against alternative topologies, a rigorous sensitivity study was performed across a parameterized continuum of stretching bounds to evaluate its impact on the metric-weighted mass matrix condition number, $\kappa(\mathbf{M})$ (Table~\ref{tab:stretch_sensitivity}).

\begin{table}[h]
\centering
\caption{Mass Matrix Conditioning and Grid Resolution Sensitivity Across Parameterized Stretching Bounds}
\label{tab:stretch_sensitivity}
\vspace{0.2cm}
\begin{tabular}{c|rrrc}
\hline
$\alpha_s$ & $\Delta z_{\min}$ (at $1.5\,\text{m}$) & $\Delta z_{\max}$ (at $z_{\max}$) & $\kappa(\mathbf{M})$ & Operational Status \\
\hline
0.50 & $12.45\,\text{m}$ & $18.20\,\text{m}$ & $1.8 \times 10^1$ & Quasi-Uniform Limit \\
1.50 & $5.12\,\text{m}$  & $11.45\,\text{m}$ & $3.4 \times 10^2$ & Weak Boundary Focus \\
\textbf{2.50} & $\mathbf{1.18\,\text{m}}$  & $\mathbf{8.92\,\text{m}}$  & $\mathbf{4.6 \times 10^3}$ & \textbf{Optimized Baseline} \\
4.00 & $0.21\,\text{m}$  & $22.40\,\text{m}$ & $8.9 \times 10^5$ & Ill-Conditioned Upper Core \\
\hline
\end{tabular}
\end{table}

The computational manifold embeds this non-uniform spacing into the discrete physics via a metric-weighted mass matrix $\mathbf{M} \in \mathbb{R}^{(N+1) \times (N+1)}$. This structure discretizes the weighted continuous inner product space using high-order Gauss--Chebyshev quadrature across $K_q = 72$ interior nodes:
\begin{equation}
M_{mn} = \int_{-1}^{1} T_m(\xi) T_n(\xi) J(\xi) \, d\xi \approx \frac{\pi}{K_q} \sum_{k=1}^{K_q} T_m(\xi_k) T_n(\xi_k) J(\xi_k)
\label{eq:mass_matrix_quad}
\end{equation}
where $T_n(\xi)$ denotes the $n$-th Chebyshev polynomial of the first kind, and $\xi_k = \cos\left(\frac{\pi(2k-1)}{2K_q}\right)$ are the roots of $T_{K_q}(\xi)$. At the optimized operational baseline, the metric-weighted system balances small-scale physical sensitivity with numerical stability ($\kappa(\mathbf{M}) = 4619$), ensuring robust Cholesky factorization ($\mathbf{M} = \mathbf{L}\mathbf{L}^T$).



\subsection{Thin SVD-Truncated Matrix Projection Architecture}
\label{sec:methods:svd}

Reconstructing high-fidelity continuous profiles from a highly constrained, sparse array of vertical observations ($M = 8$ physical tower levels) onto an $N=32$ order spectral expansion requires resolving an inherently rank-deficient linear inverse problem. Evaluating the basis functions at the fixed instrument elevations yields the rectangular evaluation design matrix $\mathbf{H} \in \mathbb{R}^{M \times (N+1)}$:
\begin{equation}
H_{ij} = T_{j-1}(\xi(z_i))
\label{eq:h_matrix_def}
\end{equation}

To minimize computational overhead while guaranteeing protection against numerical over-fitting, we reject classic normal equations or unconstrained Tikhonov diagonal dampening. Instead, we compute an economy-size, thin Singular Value Decomposition (SVD) of the rectangular operator:
\begin{equation}
\mathbf{H} = \mathbf{U}_r \mathbf{S}_r \mathbf{V}_r^T
\label{eq:thin_svd}
\end{equation}
where $\mathbf{U}_r \in \mathbb{R}^{M \times M}$ and $\mathbf{V}_r \in \mathbb{R}^{(N+1) \times M}$ retain strictly orthogonal column components, and $\mathbf{S}_r = \operatorname{diag}(s_1, s_2, \ldots, s_M)$ holds the ordered singular spectrum. This thin variant restricts matrix multiplications exclusively to the data-constrained subspace, eliminating the numerical costs associated with tracking unresolvable high-frequency spectral components.

To isolate valid physical modes from background precision noise without introducing hardcoded scaling artifacts during intense thermal modifications, we enforce a dynamic truncation tolerance criterion $\tau$:
\begin{equation}
\tau = \max\left(M, N+1\right) \cdot s_1 \cdot \epsilon_{\mathrm{mach}}
\label{eq:dynamic_tolerance}
\end{equation}
where $\epsilon_{\mathrm{mach}} \approx 2.22 \times 10^{-16}$. The effective operational rank is bounded strictly by $r_{\mathrm{eff}} = \#\{s_i > \tau\}$. The optimized, noise-filtered spectral coefficient vector $\mathbf{c}$ is computed directly via:
\begin{equation}
\mathbf{c} = \sum_{i=1}^{r_{\mathrm{eff}}} \frac{\mathbf{u}_i^T \mathbf{a}_{\mathrm{obs}}}{s_i} \mathbf{v}_i
\label{eq:svd_reconstruction}
\end{equation}
Projections are guaranteed to lie entirely within the data-validated coordinate workspace, automatically forcing unconstrained higher-order spectral modes cleanly to zero.

\subsection{Smooth Scale Partitioning and Energy Non-Conservation Caveat}
\label{sec:methods:partition}

To separate multi-scale structures without inducing non-local Gibbs ringing or artificial step discontinuities, we deploy a partition-of-unity spectral windowing framework. The global flow field is decomposed into mean synoptic ($\psi_M$), internal gravity wave ($\psi_W$), and micro-turbulent residual ($\psi_T$) scales using coupled hyperbolic tangent transition channels:
\begin{align}
  \psi_M(n) &= \frac{1}{2}\left[1 - \tanh\left(\frac{n - n_M}{\Delta}\right)\right] \label{eq:psi_m_def} \\
  \psi_W(n) &= \frac{1}{2}\left[1 + \tanh\left(\frac{n - n_M}{\Delta}\right)\right] \cdot \frac{1}{2}\left[1 - \tanh\left(\frac{n - n_W}{\Delta}\right)\right] \label{eq:psi_w_def} \\
  \psi_T(n) &= 1 - \psi_M(n) - \psi_W(n) \label{eq:psi_t_def}
\end{align}
where the scale boundaries are set to $n_M = 3$ and $n_W = 12$ with a smooth transition width $\Delta = 1.2$. By design, the partition operators sum identically to unity across the entirety of the discrete spectrum:
\begin{equation}
\psi_M(n) + \psi_W(n) + \psi_T(n) = 1 \quad \forall n \in [0, N]
\label{eq:partition_unity}
\end{equation}

Crucially, because these smooth partition vectors are applied as coefficient multipliers rather than projections onto true orthogonal eigenvectors of the mass matrix $\mathbf{M}$, \emph{global physical energy content is not strictly additive across the partitioned sub-domains}. When evaluating the individual scale components ($\mathbf{c}^{(M)} = \operatorname{diag}(\psi_M)\mathbf{c}$, $\mathbf{c}^{(W)} = \operatorname{diag}(\psi_W)\mathbf{c}$, and $\mathbf{c}^{(T)} = \operatorname{diag}(\psi_T)\mathbf{c}$), the cross-product metrics do not vanish identically:
\begin{equation}
\langle \mathbf{c}, \mathbf{M} \mathbf{c} \rangle = \sum_{j \in \{M,W,T\}} \langle \mathbf{c}^{(j)}, \mathbf{M} \mathbf{c}^{(j)} \rangle + \sum_{j \neq k} \langle \mathbf{c}^{(j)}, \mathbf{M} \mathbf{c}^{(k)} \rangle
\label{eq:energy_overlap}
\end{equation}
The off-diagonal cross-terms $\langle \mathbf{c}^{(j)}, \mathbf{M} \mathbf{c}^{(k)} \rangle \neq 0$ physically represent spatial and spectral scale overlap regions within the boundary layer transitions. While this cross-term behavior accurately mirrors the coupled wave-turbulence interaction concepts emphasized by \citet{Sun2015}, total energy balance estimates derived via this windowing scheme must explicitly track these non-orthogonal interaction components rather than assuming direct linear superposition.

\subsection{Information-Theoretic Diagnostic Features}
\label{sec:methods:diagnostics}

\subsubsection{Effective Modal Dimension ($D_{\mathrm{eff}}$) with Minimum Energy Filtering}

We compute the effective modal dimension $D_{\mathrm{eff}}$ as the exponential of the spectral Shannon entropy, converting a raw information metric into an interpretable value representing the number of actively participating, equiprobable spectral modes:
\begin{equation}
D_{\mathrm{eff}} = \exp\left( -\sum_{n=0}^{N} p_n \log p_n \right)
\label{eq:d_eff_def}
\end{equation}
where $p_n = c_n^2 / \sum_m c_m^2$ is the normalized energy in mode $n$, applying the standard information-theoretic convention that $p_n \log p_n \to 0$ when $p_n \to 0$.

Because $D_{\mathrm{eff}}$ operates as a scale-invariant shape diagnostic, a quiescent profile under weak background conditions could mathematically register an artificially inflated modal dimension due to numerical precision variations. To prevent low-amplitude instrumental noise from triggering false high-dimensional turbulence classifications, we introduce a regularized minimum energy masking operator before computing the spectral probabilities:
\begin{equation}
c_n^2 \leftarrow \begin{cases}
c_n^2, & \text{if } \sum_{m=0}^{N} c_m^2 \ge \mathcal{E}_{\text{floor}} \\
0, & \text{otherwise}
\end{cases}
\label{eq:energy_floor}
\end{equation}
where the threshold floor is defined dynamically based on historical instrument variance limits, fixed at $\mathcal{E}_{\text{floor}} = 10^{-4} \cdot \max_{t} \left(\|\mathbf{a}_{\mathrm{obs}}\|^2\right)$. This restriction forces highly quiescent, laminar configurations cleanly down to $D_{\mathrm{eff}} \to 1.0$.

\subsubsection{Wave Energy Fraction ($F_W$) and Spectral Curvature ($\chi_N$)}

The wave energy fraction isolates metric-weighted energy in the wave band relative to total energy:
\begin{equation}
F_W = \frac{\langle \mathbf{c}^{(W)}, \mathbf{M} \mathbf{c}^{(W)} \rangle}{\langle \mathbf{c}, \mathbf{M} \mathbf{c} \rangle}
\label{eq:f_w}
\end{equation}

Spectral curvature ($\chi_N$) measures the concentration of energy in high-frequency modes relative to low-frequency modes, providing a mathematical proxy for structural profile sharpness:
\begin{equation}
\chi_N = \frac{\sum_{n=2}^{N} (n/N)^2 \, c_n^2}{\sum_{n=0}^{N} (n/N)^2 \, c_n^2}
\label{eq:chi_n}
\end{equation}